\subsection*{控制面保存哪些事实}

控制面保存 b1 当前属于哪个组、归属代次、迁移标识和导入图像摘要，也可以保存工作者选主信息、服务配置与成员列表。这些记录通常比业务账本小，修改频率也低。业务组根据已经取得的配置接收正常请求，控制面短暂不可用时，有效配置下的组内业务路径可以继续推进。

控制面中的执行者会话与购买会话拥有不同对象。购买会话把账户、动作和准入期限连接起来，关闭后停止旧身份的公共准入；执行者会话表示某个进程当前能否代表某个角色工作。支付工作者的执行者会话过期后，受理事实和原平台操作继续保存，新的执行者读取同一份待办推进恢复。

\CSeventeenFigure{ch17-07-control}{控制面发布配置，业务组核验归属，资源端核验实际写入}{fig:ch17-control}

配置系统已经删除 A 的执行资格，A 先前发出的 HTTP 请求也可能正在网络中。配置变更改变的是后续可以观察到的控制事实，请求能否修改外部资源，由该资源接收请求时的规则决定。下面先展开三种协调系统保存事实的结构，再把资格接到资源写入上。

\subsection*{ZooKeeper 的节点与会话}

ZooKeeper 把小型协调数据组织成 znode 树。节点保存数据及版本，客户端可以按预期版本执行条件更新。目录路径 \code{/routes/b1} 保存归属记录，版本比较防止两个迁移者都依据同一个旧值覆盖目录。大型业务图像放在独立存储中，znode 保存图像位置、摘要和发布记录，协调树据此承担配置仲裁~\cite{zk-396-programmer}。

执行者选主可以在 \code{/executors/b1} 下创建临时顺序节点。A 得到 \code{candidate-0000000000}，B 得到 \code{candidate-0000000001}，客户端读取候选列表后，按既定顺序找到当前角色。临时节点绑定创建它的会话，会话终结时由服务端删除。顺序编号在同一父节点的创建空间中分配，适合排列这批候选。

候选者可以只观察排在自己前面的直接前驱。B 观察 A 的节点，后来的 C 观察 B 的节点。A 退出后先唤醒 B，B 重新读取列表，确认自己已经排在最前才执行 leader 流程。设置 watch 时若前驱已经消失，B 立即重读列表并重新选择前驱。这个循环接住删除与观察之间的竞争，也避免一次退出唤醒所有候选者~\cite{zk-396-recipes}。

\CSeventeenFigure{ch17-19-predecessor}{候选者观察直接前驱，通知后重读候选集再取得角色}{fig:ch17-predecessor}

编号的作用范围随对象解释。znode 版本记录该节点的数据修改次数，zxid 标识 ZooKeeper 事务次序，顺序后缀排列同一父节点下的创建。它们的生命周期和溢出条件各不相同。资源端若要把顺序后缀用作执行资格，就要先规定节点删除重建、集群历史恢复以及计数边界的处理，才能知道两个来件如何比较。

客户端与一台服务器断开连接后，先尝试在协商的会话期限内重连。成功重连可以继续使用原会话和临时节点。集群判定会话过期后，临时节点被删除，另一候选者取得角色。断线的 A 可能尚不知道自己已过期，要等重新连接才得到相应通知。A 的本地变量及其已发出的外部请求，可以比临时节点存在更久。

顺序节点创建遇到连接错误还会产生结果不确定：服务端可能已创建节点，客户端尚未收到路径。参与者可以把自己生成并持久保存的唯一标记写入节点名或内容，重连后搜索已创建的候选，再决定是否重新创建。仅在超时后无条件创建另一候选，会让同一参与者留下多个节点，退出与接管也随之含混~\cite{zk-396-recipes}。

标准 watch 是一次性通知。B 读取 \code{/routes/b1} 并设置 watch，A 将归属从 G1 改为 G2，B 收到通知后重读当前数据并重新设置 watch。通知到达时，数据还可能已从 G2 改成 G3。B 因而把通知作为重新核读事实的触发器。断线期间若节点被创建又删除，客户端还可能看不到这个完整中间过程。

\CSeventeenFigure{ch17-08-session-watch}{会话过期改变候选资格，一次性 watch 触发重读与重新订阅}{fig:ch17-session-watch}

重新读取和重新设置 watch 使用支持这一组合的读取接口，避免先读后另行订阅之间出现空隙。支持持久 watch 的版本提供不同的通知生命周期，应用围绕当前事实重建状态。

ZooKeeper 的读可以由所连接服务器提供，跨客户端在某个时刻可能看见不同进度。A 成功把目录从 G1 改为 G2，再用另一条通信告诉 B，B 直接读取自己连接的落后服务器时可能看见 G1。官方文档给出的协调做法是先执行 sync，再在同一客户端路径上读取。系统需要建立这类跨客户端顺序时，显式执行同步步骤，并按部署版本核对读取保证~\cite{zk-396-programmer}\footnote{本节按 ZooKeeper 3.9.6 指南与 Recipes 讨论标准 watch。官方说明顺序后缀使用父节点内的计数，并列出 32 位有符号整数的溢出边界；长期资格比较应显式覆盖这一生命周期。}。

\subsection*{etcd 的事务、修订号与租约}

etcd 的 KV 接口提供条件事务 Txn。迁移者读到 b1 的旧目录，再提交一项事务：值等于旧值时写入新值，否则读取当前值。两个迁移者并发交接时，先成功者改变当前值，后一个比较就会失败。比较与更新属于同一个 etcd 事务，入口可以根据失败分支取得已经成立的新目录~\cite{etcd-37-api}。

etcd 为 KV 修改分配集群范围内的 revision。同一事务中的修改共享一个修订号，单键还保存 \code{create\_revision}、\code{mod\_revision} 和 \code{version}。首次创建锁键时，可用创建修订号标记这次锁实例，修改修订号描述该键最近一次变化的进度。\code{lease ID} 引用租约，本身没有有序代次的含义。

默认 Range 走线性化读取路径，反映集群已经成立的 KV 次序。显式选择 \code{serializable} 的成员本地读取，可以提高本地读取可用性，但可能落后于当前多数。目录缓存允许落后，业务组会核验归属代次；迁移者决定目标激活前，则沿线性化路径确认当前目录，把发布依据与业务资格连接起来~\cite{etcd-37-guarantees}。

租约把若干键的生命周期绑定到执行者的存活维护。A 取得租约，在条件事务中首次创建执行者键并附着该租约。A 续租时键可继续存在，租约到期或被撤销后关联键被删除。B 随后创建新执行者键，取得另一创建修订号。业务分区的 $S$ 绑定和原付款结果独立保存，租约回收只回收其所关联的控制键。

\subsection*{watch 缓存的完整重建}

目录观察者先执行一次线性化 Range，取得完整目录及响应修订号 $r$，再从 $r+1$ 建立 watch。两步之间发生更新时，watch 可以从保留的历史中补齐。观察者按完整修订号应用事件，保存最近完整处理位置。重连从该位置之后订阅，避免只处理了同一事务的一部分就跳过其余事件。

\CSeventeenFigure{ch17-09-watch-rebuild}{Range 取得基线，watch 衔接后缀，压缩后重新建立完整基线}{fig:ch17-watch-rebuild}

watch 有自己的交付进度。目录写入已经完成，观察者收到对应事件的时间却没有固定上界。应用利用进度通知判断已经覆盖的修订号，暂时没有事件只能说明当前尚未收到新通知。作交接决定的管理路径使用合适的 KV 读取，普通缓存则按其刷新规则继续接收或重建。

历史会被压缩。观察者离线较久，恢复所用 \code{start\_revision} 早于可用窗口时，etcd 取消 watch 并返回 \code{compact\_revision}。观察者重新执行 Range，替换完整目录，再从新响应修订号之后建立 watch。反复请求已压缩的位置，只会再次得到取消。

分批读取大目录时，所有分页要固定在同一修订号。第一页取得 $r$ 后，后续页显式读取 $r$ 对应的视图，完整图像完成后再从 $r+1$ 衔接事件。若每一页都读各自最新状态，删除、搬迁和新增可能交错成一份从未存在过的目录。

完整基线与完整修订衔接后，观察者可以发布带水位的目录缓存。迁移者据目录执行交接决定时，则直接沿线性化读取核验归属、代次与迁移记录，避免把尚未送达的watch事件误判为目录没有变化~\cite{etcd-37-api,etcd-37-guarantees}\footnote{etcd 的 KV 操作和 watch 使用不同的保证描述。watch 可续接保留的历史，也可能任意延迟，其保证不等于在事件交付时进行一次线性化读取。}。

\subsection*{成员变化如何进入控制面的多数}

协调集群自身也会换机器。尚未同步的新节点立即进入投票集合，会提高多数门槛，却没有立即增加能响应的有效副本。etcd 提供 learner 角色，让新节点先接收日志和状态，追上之后再提升为投票成员。提升由服务端检查条件，未同步时会被拒绝~\cite{etcd-37-reconfigure}。

三投票成员集群增加一个 learner 时，原来的三成员多数继续工作。新成员追赶并满足条件后提升为第四个投票成员，多数门槛随已提交配置改变。移除也作为管理操作执行。管理者按产品规则串行变更，确认前一项已经成立再开始下一项。

\CSeventeenFigure{ch17-20-learner}{新成员先追赶，再通过资格检查进入投票集合}{fig:ch17-learner}

成员列表和业务目录各自保存自己的事实。增加 etcd 成员会多复制一份相同目录，b1 的业务状态保持在原购买组。购买组增加 Raft 副本，则改变该组的承载与投票集合，目录中的 b1 归属可以保持原值。运维据此区分控制集群成员变更和业务分区搬迁。

\subsection*{Consul 的 session 与读选项}

Consul 的 session 可以关联节点、健康检查和 TTL，再通过 KV 的 acquire 取得建议性锁。成功后，键记录持有 session 及 LockIndex。相同 session 再次 acquire 可更新数据，释放或销毁 session 则按配置清除关联，后继参与者再 acquire 取得自己的锁实例。TTL 是过期判定下界，实际清理可以更晚~\cite{consul-sessions}。

健康检查帮助故障后的接管，也可能把存活但暂时失联的进程判为失效。lock-delay 在旧 session 失效后延迟重新取得锁，给旧持有者观察失效并停止工作的机会。延迟缓和时间重叠，外部 HTTP 到达后的效果由资源端规则决定。

acquire 建立协作参与者遵守的锁约定。普通客户端即使没持有 session，也可以直接写入或删除同一 KV 键。资源若使用 \code{(Key, LockIndex, Session)} 识别锁实例，就要保护键的完整生命周期和修改权限，并核对来件所携带的实例。删除重建、普通写入或恢复集群历史，都可能越过只比较 LockIndex 的使用前提。

Consul 读取分为 default、consistent 和 stale。consistent 由 leader 向多数确认领导资格后处理，增加相应通信；stale 允许任意服务器本地回答，陈旧程度没有固定上界，失去 leader 时也可能读到旧值。default 存在旧 leader 尚未发现失去资格的窗口。迁移发布判断和近似服务发现根据读取用途选择不同路径~\cite{consul-consistency}。

三种系统都能保存协调事实，节点或键的生命周期、通知和读取条件各有具体规则。配置适配器据此实现读取、条件更新、重连和重建；业务步骤明确记录自己已得到哪一种事实，才能在结果不确定或会话过期后选择下一项动作。

\subsection*{把资格落实到资源端}

设 A 持有代次 17，发出库存请求后暂停。控制面撤销 A 的资格，B 取得代次 23。B 先在资源端安装 23，再提交库存修改。A 的旧 HTTP 到达时，资源将来件中的 17 与已保存的 23 比较，在修改前返回 \code{stale\_fence}。这一步称为 fencing，即由被保护资源阻挡已经失去资格的写入。

\CSeventeenFigure{ch17-10-resource-fence}{新资格先安装到资源，迟到的旧请求在修改前被拒绝}{fig:ch17-resource-fence}

资源端把资格比较和实际修改放在同一原子边界。若先检查 17 有效，释放锁后再改库存，B 可以在两步之间安装 23，旧 A 随后依然完成修改。下面给出单个库存资源 \code{inventory-1} 的事务顺序，绑定表属于这个资源的保存域。入口已核验身份、输入类型与范围，\code{requested\_quantity} 为正整数；\code{S\_key} 是四分量身份的规范表示，\code{input\_bytes} 是需要绑定的规范输入。

\begin{BookCode}
BEGIN;
lock resource_row('inventory-1');
if request.fence != resource_row.current_fence:
    ROLLBACK; return stale_fence;
old = find binding by S_key;
if old exists:
    COMMIT;
    return old.result if old.input_bytes == input_bytes else conflict;
if stock < requested_quantity:
    result = {status: rejected, reason: inventory_unavailable};
    insert binding(S_key, input_bytes, result);
    COMMIT; return result;
decrease stock by requested_quantity;
effect_id = insert effect and obtain its stored identifier;
result = {effect_id, original_S, original_input, status: committed};
insert binding(S_key, input_bytes, result);
COMMIT;
return result;
\end{BookCode}

安装新代次的管理操作在同一资源行上序列化，只允许向前推进。\code{current\_fence} 决定资源端的资格接收顺序，业务绑定保存同 $S$ 的原输入和原结果。B 重试时先通过资格核验，再读取原结果，库存只减少一次。A 的旧代次来件在修改前退出。业务拒绝也保存为原结果，同一身份的重试便能解释此前的库存判断。

这里分别检查执行资格和业务身份。资格代次限制谁能修改资源，$S$ 绑定决定同一业务如何恢复。A 与 B 若都按同一 $S$ 恢复付款，可以共同访问平台的原幂等入口；B 若换了业务身份，即使资格最新，也可能合法创建第二项业务。平台只有业务幂等键而没有资源 fence 接口时，接管者根据该平台的实际能力，沿原键和原证据推进。

\CSeventeenFigure{ch17-21-identity-qualification}{归属、工作者资格与平台身份在各自接收点检查}{fig:ch17-identity-qualification}

代次比较限定在来源的连续历史中。一个 etcd 集群内的创建修订号可以区分先后锁实例，恢复为旧集群或新集群后，裸修订号可能重新落入资源见过的范围。灾难恢复从不可回退来源取得更高的集群代次，与局部修订号组成可排序资格，或者先冻结资源再受控交接。资源保存的最高资格也属于恢复责任，恢复者同时核对资格来源与资源水位，随后开放修改。
\FloatBarrier
